\printedpage{14}
% head: A. Grothendieck : Dérivateurs, Ch. XII, Édition des Universités de Montpellier II et Paris VII
least expected: the fibration axiom. Let us take the weak form: Let $f : X \longrightarrow Y$ be a $W_\omega$-fibration and $f \in W_i$, are the fibers $X_s$ $W_i$-aspheric (\footnote[37]{$\omega > \infty$, $W_\omega \subset W_\infty \subset W_n$})? We may suppose $X$ and $Y$ 0-connected, let $x \in X_s$, the homotopy exact sequence gives
\[
\pi_{i+1}(X) \xrightarrow[\text{if } i \leq n-1]{\text{iso}} \pi_{i+1}(Y) \longrightarrow \pi_i(X_s) \longrightarrow \pi_i(X) \xrightarrow[\text{if } i \leq n]{\text{iso}} \pi_i(Y),
\]
so
\[
\pi_i(X_s) = 1 \qquad \text{if } 0 \leq i \leq n-1,
\]
but for $i = n$ we find\msp{page 17}
\[
\pi_n(X_s) = \mathrm{Coker}\left(\pi_{n+1}(X) \longrightarrow \pi_{n+1}(Y)\right),
\]
there is no reason for this to be zero, as it should if $X_s \longrightarrow e$ were in $W_n$.

It is true that sometimes one does strengthen the condition for $W_n$, let us say one defines $W'_n \subseteq W_n$ by
\[
f \in W'_n \iff \left\{
\begin{array}{ll}
\text{a)} & \pi_0(f) \text{ bijective,} \\[1ex]
\text{b)} & \pi_i(f,x) : \pi_i(X,x) \longrightarrow \pi_i(Y,f(x)) \text{ bijective for } 0 < i \leq n \text{ and } x \in X, \\[1ex]
\text{c)} & \pi_{n+1}(f,x) : \pi_{n+1}(X,x) \longrightarrow \pi_{n+1}(Y,f(x)) \ \textit{surjective,}
\end{array}
\right.
\]
that is to say also
\[
f \in W'_n \iff \left\{
\begin{array}{l}
\text{the ``homotopy fibers'' } Z \text{ of } X \longrightarrow Y \text{ (i.e. the fibers of } \bar{X} \xrightarrow{\bar{f}} Y\text{)} \\
\text{are } n\text{-connected, i.e. 0-connected and } \pi_i(Z_y,z) = 0 \text{ for } 1 \leq i \leq n.
\end{array}
\right.
\]
Then $X$ in Cat will be $W'_n$-aspheric if and only if it is $W_n$-aspheric, i.e. $n$-connected. So by form 2 of the definition of $W'_n$, axiom W(4b$''$) is satisfied. It must go wrong elsewhere! Perhaps saturation already.

\hspace*{3em}If $gf$ and $g$ are in $W'_i$, the same is not necessarily true of $f$,

for if $\bar{g}\bar{f}$ (and hence $\bar{g}$) is surjective, the same is not necessarily true for $\bar{f}$! So $W'_n$ does not deserve the name of ``localizer''. % [sic] par forcément -> pas forcément, misprint

On the other hand, $W_n$ seems excellent to me, if one does not ask the impossible of it - to satisfy a ``fibration axiom'' which does seem (more or less) to characterize $W_\infty$ alone among the ``good'' localizers. In any case the connectedness axiom W(5) is satisfied, and the filtering limits axiom W(6) as well. Moreover W(1) (and even strong saturation), W(2), W(3) are satisfied by every $W_{\mathbf{D}}$ associated with a derivator - which confirms their importance once again. The axioms W(4) (the ``draconian'' fibration axiom, $\pm$ special to $W_\infty$), W(5) (connectedness) and % [sic] connéxité -> connexité, misspelling (written connexité elsewhere)
\printedpage{15}
% head: A. Grothendieck : Dérivateurs, Ch. XII, Édition des Universités de Montpellier II et Paris VII
\msp{page 18}W(6) (limits) appear to me to be of a more special nature, I would not think of imposing any of them for the terminological use of ``basic localizer''.

\newpage
\printedpage{16}
% head: A. Grothendieck : Dérivateurs, Ch. XII, Édition des Universités de Montpellier II et Paris VII
\section*{II\quad Perplexities about the axioms and the terminology.}

In I and most of the time, in working with derivators, I have assumed \fbox{$\ W_{\mathbf{D}} \supseteq W_\infty\ $}, a hypothesis which seemed essential to me (wrongly (\footnote[38]{\textit{wrongly} it now turns out.}) or rightly). So we see that if $W_{\mathbf{D}}$ satisfies the fibration axiom W(4), even in its weakest form W(4b$''$), it is \textit{equal} to $W_\infty$. This does show that this axiom must be seen as practically a \textit{characterization} of $W_\infty$ (among the basic localizers which contain $W_\infty$ -- and these localizers seem to me the most important, and perhaps the only really useful ones).

The theory of \textbf{D}-fibrations in I, p. 111-126, was done in the spirit just mentioned, assuming $W_{\mathbf{D}} \supseteq W_\infty$. The notion of \textbf{D}-fibrations on which I had finally settled was consequently different from the notions considered here in a different spirit (page 4). It amounts to testing condition a) (almost always too strong, except for $W = W_\infty$), characterizing ultrastrong $W$-fibrations, by $Y'' \longrightarrow Y'$, where $Y''$, $Y'$ are both $W_\infty$-aspheric (which is more demanding than $W$-aspheric). It is only when $Y''$ and $Y'$ are $W_\infty$-aspheric (instead of $W$-aspheric), that one asks that $X'' \longrightarrow X'$ be in $W$. We thus find\msp{page 19} a less demanding notion than that of strong $W$-fibration. If we denote this condition here by b), we find (provided that $W_\infty \subseteq W$), with the notations of page 4 (\footnote[39]{Cf. the right notion of $W$-fibration in a context where one does not assume $W_\infty \subseteq W$, in th. 3 and def. 3, p. 51, 52.}) % [sic] arreté -> arrêté, missing circumflex
% In the print, a curved arrow starts at "b" in the diagram below and points up to "b)" in the line "par b), on trouve ..." above; it is drawn here to the empty node at the top left.
\[
\begin{tikzcd}[column sep=small, row sep=large]
{} & & & & \\
& & \mathrm{a} \arrow[d, Rightarrow] & & W\text{-ultrastrong fibration} \\
& & \mathrm{a}' \arrow[d, Rightarrow] & & W\text{-strong fibration} \\
& & \mathrm{b} \arrow[dl, Rightarrow] \arrow[dr, Rightarrow] \arrow[uuull, bend left=30] & & W\text{-fibration} \\
W\text{-left fibration} & \mathrm{a}_1 & & \mathrm{a}_2 & W\text{-right fibration.}
\end{tikzcd}
\]
a$_1$), a$_2$) were also called cohomological \textbf{D}-fibrations resp. homological \textbf{D}-fibrations (\footnote[40]{No, this is a still stronger notion, by dint of locally constant coefficients, cf. I, p. 121.}). When \textbf{D} satisfies Der 3a$'$ (I, p. 111) on locally constant \textbf{D}-coefficients on the $X$ that are $W_\infty$-aspheric, then the notions b), a$_1$), a$_2$) are equivalent. (I presume this is the case for all the \textbf{D} and all the $W$ that are really interesting.) On the other hand, I have the impression that the notion of strong $W$-fibration, or even ultrastrong, is hardly usable except in the case $W = W_\infty$, where it coincides with the preceding ones. That is to say that one can probably forget the terminology, and remember on the other hand that in the case $W = W_\infty$, the % [sic] cohomologique -> cohomologiques, agreement with "D-fibrations"
\printedpage{17}
% head: A. Grothendieck : Dérivateurs, Ch. XII, Édition des Universités de Montpellier II et Paris VII
$W$-fibrations have properties of extreme strength, expressed by the two properties singled out in W(4a) and W(4b).

\msp{page 20}In the theory of $\mathrm{HOT}_W$, influenced by the general reflection on derivators and axiom Der 3b, I have worked under the hypothesis $W \supseteq W_\infty$, but using a fibration axiom L5 $=$ W(4b) which, it turns out, ultimately implies $W = W_\infty$. I have the impression that so far (properties of $f_!$ for HOT) I have not used $W \supseteq W_\infty$, so it may be useful, since ``it costs no more'' (and forces one to be more attentive to what one uses), to drop $W \supseteq W_\infty$ (implying $W = W_\infty$), and on the other hand to put on $W$ axioms as strong (notably of fibration) as the properties of $W_\infty$, ``without depriving oneself of anything''. So I shall assume W(1), W(2), W(3), W(4a) (the strongest form of the fibration axiom, cf. pages 13, 14) - that is all that has been useful up to now. If one adds W(5) (connectedness axiom) to this it implies $W \subseteq W_\infty$. But for the moment I do not think I have a need for it.

But meditating on the question, I believe I have found the

\textbf{Theorem 2} (!) (\footnote[41]{Cf. for a stronger result XV 9.12, p. 133 bis. \textbf{NB} The result is more general, but not really stronger, since it does not give $W = \Fl(\Cat)$.}): \textit{Let} $W \subseteq \Fl(\Cat)$ \textit{satisfy} W(1,2,3,4b$''$)\textit{. If} $W$ \textit{does not satisfy} W(5)\textit{, then} $W = \Fl(\Cat)$.

\textsc{Proof}. Let $f : X \longrightarrow Y$ be in Cat, such that $\pi_0(f)$ is not bijective. Let $I$ be the discrete category having as set of objects the set $\pi_0(Y)$, so we have a commutative diagram
\[
\begin{tikzcd}[column sep=small]
X \arrow[rr, "f"] \arrow[dr, "p"'] & & Y \arrow[dl, "q"] \\
& I, &
\end{tikzcd}
\]
on the other hand a functor from a category into a discrete category is always $W_\omega$-perfect, so $p$, $q$ are $W$-fibrations. Since\msp{page 21} $f \in W$, it follows that $f_i \in W$ for all $i \in I$ (\footnote[42]{We use W(4) in a form stronger than W(4b$''$)!}).

If $\pi_0(f)$ were not surjective, taking $i$ which is not in the image, we would have $X_i = \emptyset \longrightarrow Y_i \neq \emptyset$, which would be in $W$. So let
\[
f : X = \emptyset \longrightarrow Y
\]
be in $W$. Then $\forall$ $Z$ in Cat, $f \times \mathrm{id}_Z : \underbrace{X \times Z}_{= \emptyset} \longrightarrow Y \times Z$ is in $W$ (since $Z$ is
\printedpage{18}
% head: A. Grothendieck : Dérivateurs, Ch. XII, Édition des Universités de Montpellier II et Paris VII
$W$-fibered over $e$), so $\emptyset \longrightarrow Y \times Z$ is in $W$. So too the projection
$\begin{tikzcd}[column sep=small, row sep=small, baseline=(\tikzcdmatrixname-1-1.base)]
Y \times Z \arrow[rr] & & Y. \\
& \emptyset \arrow[ul] \arrow[ur] &
\end{tikzcd}$
Since this projection is a $W$-fibration, it follows by W(4b$''$) that its fibers are $W$-aspheric. Since $Y \neq \emptyset$, there is such a fiber, it is $\simeq Z$. So $Z$ is $W$-aspheric. So we find that \textit{every} category $Z$ in Cat is $W$-aspheric. So every arrow $Z \longrightarrow Z'$ in Cat is in $W$, q.e.d.

Suppose now $\pi_0(f)$ surjective, but not injective. There would be a fiber of $\pi_0(X) \longrightarrow \pi_0(Y) = I$ which would be of cardinality $> 1$. Taking the corresponding $f_i$, we find $X_i \longrightarrow Y_i$ with $Y_i$ 0-connected and $\card \pi_0(X_i) \geq 2$. So we are reduced to the case
\[
f : X \longrightarrow Y \qquad\qquad \left\{
\begin{array}{l}
f \in W \\[1ex]
Y \text{ 0-connected} \\[1ex]
\card \pi_0(X) \geq 2.
\end{array}
\right.
\]
Let us consider the canonical factorization
\[
X \xrightarrow{i_f} \bar{X} \xrightarrow{\bar{f}} Y,
\]
with $i_f \in W$, so $\bar{f} \in W$. So by W(4b$''$) the fibers of $\bar{f}$ are $W$-aspheric. But we see at once, on decomposing $\bar{X}$ according to the connected components $X_i$ of $X$, that these fibers have a $\pi_0$ of cardinality $\geq \card \pi_0(X) \geq 2$. So we are reduced to the case where $Y = e$, i.e.
\[
X \ W\text{-aspheric}, \qquad \card \pi_0(X) \geq 2.
\]
Let us consider $X \times X$, which is also $W$-aspheric, and the diagonal map\msp{page 22}
\[
d : X \longrightarrow X \times X = Y,
\]
which is therefore in $W$. But $\pi_0(d)$ is identified with the diagonal map $\pi_0(X) \longrightarrow \pi_0(X) \times \pi_0(X)$, which is not surjective since $\card \pi_0(X) \geq 2$. So for the first part, $W = \Fl(\Cat)$.

So in more positive form:

\textbf{Corollary.} \textit{Let} $W \subseteq \Fl(\Cat)$ \textit{satisfy} W(1,2,3,4b$''$)\textit{, and be such that} $W \neq \Fl(\Cat)$. \textit{Then $W$ satisfies the connectedness axiom} W(5)\textit{, and consequently (th. 1, page 12)} $W \subseteq W_\infty$.